\ifdefined\gkver\else\def\gkver{student}\fi
\documentclass[\gkver]{gaokaozhenti}
\begin{document}

\chapter{2012年北京理综}
%% paperType: 理综物理部分

\begin{enumerate}
\item
%% number: 13
%% paperName: 2012年北京理综第 13 题 6 分
%% typeId: 1
%% type: 单选题
%% chapter: 原子和原子核
%% point: 玻尔模型
%% method: 无
%% score: 6
%% degree: 850
%% duplicateId: 0
%% body:
一个氢原子从 $n=$3 能级跃迁到 $n=$2 能级，该氢原子\xzanswer{B}
\fourchoices[answer=B]
{放出光子，能量增加}
{放出光子，能量减少}
{吸收光子，能量增加}
{吸收光子，能量减少}
%% answer: B
%% memo:
\memoanswer{
原子由高能级 3 跃迁到低能级 2 的过程中，原子能量减少，放出光子。答案 B。
}

\item
%% number: 14
%% paperName: 2012年北京理综第 14 题 6 分
%% typeId: 1
%% type: 单选题
%% chapter: 光学
%% point: 光的电磁说
%% method: 无
%% score: 6
%% degree: 850
%% duplicateId: 0
%% body:
一束单色光经由空气射入玻璃，这束光的\xzanswer{A}
\fourchoices[answer=A]
{速度变慢，波长变短}
{速度不变，波长变短}
{频率增高，波长变长}
{频率不变，波长变长}
%% answer: A
%% memo:
\memoanswer{
单色光由光疏介质空气进入光密介质玻璃，频率不变，但介质对光的折射率增大，由 $n=\frac{c}{v}$ 及 $v=\lambda f$，可知光的波长和速度都减小，答案 A。
}

\item
%% number: 15
%% paperName: 2012年北京理综第 15 题 6 分
%% typeId: 1
%% type: 单选题
%% chapter: 交流电
%% point: 交流电
%% method: 无
%% score: 6
%% degree: 650
%% duplicateId: 0
%% body:
一个小型电热器若接在输出电压为 $10\UV$ 的直流电源上，消耗电功率为 $P$；若把它接在某个正弦交流电源上，其消耗的电功率为 $\frac{P}{2}$。如果电热器电阻不变，则此交流电源输出电压的最大值为\xzanswer{C}
\fourchoices[answer=C]
{$5\UV$}
{$5\sqrt{2}\UV$}
{$10\UV$}
{$10\sqrt{2}\UV$}
%% answer: C
%% memo:
\memoanswer{
电热器电阻不变，根据功率表达式和交流电的有效值得：$P=\frac{(10\UV)^{2}}{R}$；$\frac{P}{2}=\frac{(U_{m}/\sqrt{2})^{2}}{R}$，可求得 $U_{m}=10\UV$，答案 C。
}

\item
%% number: 16
%% paperName: 2012年北京理综第 16 题 6 分
%% typeId: 1
%% type: 单选题
%% chapter: 磁场
%% point: 带电粒子在磁场中的运动
%% method: 等效替代
%% score: 6
%% degree: 650
%% duplicateId: 0
%% body:
处于匀强磁场中的一个带电粒子，仅在磁场力作用下做匀速圆周运动。将该粒子的运动等效为环形电流，那么此电流值\xzanswer{D}
\fourchoices[answer=D]
{与粒子电荷量成正比}
{与粒子速率成正比}
{与粒子质量成正比}
{与磁感应强度成正比}
%% answer: D
%% memo:
\memoanswer{
粒子的运动可等效为环形电流，粒子在一个周期内只通过某一个截面一次，则环形电流在一个周期 $T$ 内的电量为 $q$，由电流定义式 $I=\frac{q}{T}$。粒子在洛伦兹力作用下做匀速圆周运动，有 $qvB=\frac{mv^{2}}{r}$，$T=\frac{2\pi r}{v}$，联立解得 $I=\frac{Bq^{2}}{2\pi m}$。环形电流与磁感应强度成正比，与粒子质量成反比，与粒子电荷量的平方成正比，而与粒子速率无关，答案 D。
}

\item
%% number: 17
%% paperName: 2012年北京理综第 17 题 6 分
%% typeId: 1
%% type: 单选题
%% chapter: 机械振动
%% point: 振动图象
%% method: 无
%% score: 6
%% degree: 650
%% duplicateId: 0
%% body:
一个弹簧振子沿 $x$ 轴做简谐运动，取平衡位置 O 为 $x$ 轴坐标原点。从某时刻开始计时，经过四分之一的周期，振子具有沿 $x$ 轴正方向的最大加速度。能正确反映振子位移 $x$ 与时间 $t$ 关系的图像是\xzanswer{A}
\fourchoices[answer=A, ispicture=true, h=2.2cm]
{figs/17a.png}
{figs/17b.png}
{figs/17c.png}
{figs/17d.png}
%% answer: A
%% memo:
\memoanswer{
由简谐运动知 $a = -\frac{kx}{m}$，则在 $t = \frac{T}{4}$ 时刻，振子具有沿 $x$ 轴正方向的最大加速度（正的最大），它的位移为沿 $x$ 轴负方向的最大位移（负的最大），满足条件的图像只有 A，答案 A。
}

\item
%% number: 18
%% paperName: 2012年北京理综第 18 题 6 分
%% typeId: 1
%% type: 单选题
%% chapter: 曲线运动
%% point: 人造地球卫星
%% method: 无
%% score: 6
%% degree: 650
%% duplicateId: 0
%% body:
关于环绕地球卫星的运动，下列说法正确的是\xzanswer{B}
\fourchoices[answer=B]
{分别沿圆轨道和椭圆轨道运行的两颗卫星，不可能具有相同的周期}
{沿椭圆轨道运行的一颗卫星，在轨道不同位置可能具有相同的速率}
{在赤道上空运行的两颗地球同步卫星，它们的轨道半径有可能不同}
{沿不同轨道经过北京上空的两颗卫星，它们的轨道平面一定会重合}
%% answer: B
%% memo:
\memoanswer{
所有的同步卫星都在同一个赤道轨道上运动，C 错误；沿不同轨道经过北京上空的两颗卫星，它们的运行轨道面与赤道面的夹角可以不同，它们的轨道平面就不会重合，D 错误；分别沿圆轨道和椭圆轨道运行的两颗卫星，可能具有相同的周期，A 错误；沿椭圆轨道运行的一颗卫星，在轨道的关于长轴对称的两个位置的速率相等，所以在轨道不同位置可能具有相同的速率是正确的。答案 B。
}

\item
%% number: 19
%% paperName: 2012年北京理综第 19 题 6 分
%% typeId: 1
%% type: 单选题
%% chapter: 电磁感应
%% point: 楞次定律推广
%% method: 无
%% score: 6
%% degree: 650
%% duplicateId: 0
%% body:
物理课上，老师做了一个奇妙的“跳环实验”。如图，她把一个带铁芯的线圈 $L$、开关 $S$ 和电源用导线连接起来后，将一金属套环置于线圈 $L$ 上，且使铁芯穿过套环，闭合开关 $S$ 的瞬间，套环立刻跳起。某同学另找来器材再探究此实验。他连接好电路，经重复实验，线圈上的套环均未动，对比老师演示的实验，下列四个选项中，导致套环未动的原因可能是\xzanswer{D}
\onepicture[h=3.0cm, align=right]{figs/19.png}
\fourchoices[answer=D]
{线圈接在了直流电源上}
{电源电压过高}
{所选线圈的匝数过多}
{所用套环的材料与老师的不同}
%% answer: D
%% memo:
\memoanswer{
在开关闭合的瞬间，线圈中的电流变大，磁场变强，穿过套环的磁通量变大，在金属套环内产生感应电流。感应磁场必然阻碍原磁场的增大，所以金属套环会受到线圈的斥力而跳起。在实验时一般采用直流电源，电压不能太大（以不烧导线和电源的条件下现象明显），所选线圈的匝数越多，现象也越明显。如果该学生所用套环的材料为非金属，则不会观察到“跳环实验”。答案 D。
}

\item
%% number: 20
%% paperName: 2012年北京理综第 20 题 6 分
%% typeId: 1
%% type: 单选题
%% chapter: 其他
%% point: 物理思想
%% method: 无
%% score: 6
%% degree: 450
%% duplicateId: 0
%% body:
“约瑟夫森结”由超导体和绝缘体制成．若在结两端加恒定电压 $U$，则它会辐射频率为 $\nu$ 的电磁波，且 $\nu$ 与 $U$ 成正比，即 $\nu=kU$。已知比例系数 $k$ 仅与元电荷的 2 倍和普朗克常数 $h$ 有关，你可能不了解此现象的机理，但仍可运用物理学中常用的方法，在下列选项中，推理比例系数的值可能为\xzanswer{B}
\fourchoices[answer=B]
{$\frac{h}{2e}$}
{$\frac{2e}{h}$}
{$2he$}
{$\frac{1}{2he}$}
%% answer: B
%% memo:
\memoanswer{
物理公式表达了各物理量间的质量和单位双重关系，光子的能量与光的频率成正比 $E = h\nu$，电场力对电子所做电功为 $W = eU$，又 $\nu = kU$，所以 $k = \frac{\nu}{U} = \frac{E / h}{W / e}$。由于 $E$ 和 $W$ 有相同的单位，所以 $k$ 的单位与 $\frac{e}{h}$ 的单位相同。根据题意 $k$ 仅与元电荷 $e$ 的 2 倍和普朗克常量 $h$ 有关，答案 B。
}

\item
%% number: 21
%% paperName: 2012年北京理综第 21 题 18 分
%% typeId: 4
%% type: 实验题
%% chapter: 电路
%% point: 电路实验
%% method: 无
%% score: 18
%% degree: 450
%% duplicateId: 0
%% body:
在“测定金属的电阻率”实验中，所用测量仪器均已校准．待测金属丝接入电路部分的长度约为 $50\Ucm$。
\begin{enumerate}
	\item
	用螺旋测微器测量金属丝的直径，其中某一次测量结果如图 1 所示，其读数应为\_\_\_\_\_\_\_\_\_\_\_mm（该值接近多次测量的平均值）。
	\onepicture[h=2.8cm, align=right]{figs/21a.png}

	\item
	用伏安法测金属丝的电阻 $R_{x}$。实验所用器材为：电池组（电动势为 $3\UV$，内阻约 $1\UO$）、电流表（内阻约 $0.1\UO$）、电压表（内阻约 $3\UkO$）、滑动变阻器 $R$（0\textasciitilde{}20$\UO$，额定电流 $2\UA$）、开关、导线若干。

	某小组同学利用以上器材正确连接好电路，进行实验测量，记录数据如下：

	\begin{center}
	\begin{tabular}{|c|c|c|c|c|c|c|c|}
	\hline
	次数 & 1 & 2 & 3 & 4 & 5 & 6 & 7 \\
	\hline
	$U$/V & 0.10 & 0.30 & 0.70 & 1.00 & 1.50 & 1.70 & 2.30 \\
	\hline
	$I$/A & 0.020 & 0.060 & 0.160 & 0.220 & 0.340 & 0.460 & 0.520 \\
	\hline
	\end{tabular}
	\end{center}

	由以上数据可知，他们测量 $R_{x}$ 是采用图 2 中的\_\_\_\_\_\_\_\_\_\_\_\_图（选填“甲”或“乙”）。
	\onepicture[h=2.2cm, align=right]{figs/21b.png}

	\item
	图 3 是测量 $R_{x}$ 的实验器材实物图，图中已连接了部分导线，滑动变阻器的滑片 $P$ 置于变阻器的一端。请根据图（2）所选的电路图，补充完成图 3 中实物间的连线，并使闭合开关的瞬间，电压表或电流表不至于被烧坏。
	\onepicture[h=2.8cm, align=right]{figs/21c.png}

	\item
	这个小组的同学在坐标纸上建立 $U$、$I$ 坐标系，如图 4 所示，图中已标出了测量数据对应的 4 个坐标点。请在图 4 中标出第 2、4、6 次测量数据坐标点，并描绘出 $U-I$ 图线。由图线得到金属丝的阻值 $R_{x}=$\_\_\_\_\_\_\_\_\_\_\_$\Omega$（保留两位有效数字）。
	\onepicture[h=2.4cm, align=right]{figs/21d.png}

	\item
	根据以上数据可以估算出金属丝的电阻率约为\_\_\_\_\_\_\_\_\_\_\_（填选项前的符号）。

	（A）$1\times10^{-2}\UOm$\quad（B）$1\times10^{-3}\UOm$\quad（C）$1\times10^{-6}\UOm$\quad（D）$1\times10^{-8}\UOm$

	\item
	任何实验测量都存在误差．本实验所用测量仪器均已校准，下列关于误差的说法中正确的选项是\_\_\_\_\_\_\_\_\_\_\_（有多个正确选项）。

	（A）用螺旋测微器测量金属丝直径时，由于读数引起的误差属于系统误差

	（B）由于电流表和电压表内阻引起的误差属于偶然误差

	（C）若将电流表和电压表内阻计算在内，可以消除由测量仪表引起的系统误差

	（D）用 $U-I$ 图像处理数据求金属丝电阻可以减小偶然误差
\end{enumerate}
%% answer: 
\jdanswer{
\begin{enumerate}
	\item
	0.395 \textasciitilde{} 0.399

	\item
	甲

	\item
	如答图 3
	\onepicture[h=2.8cm, align=right]{figs/21an1.png}

	\item
	如答图 4（4.3 \textasciitilde{} 4.7）
	\onepicture[h=2.4cm, align=right]{figs/21an2.png}

	\item
	C

	\item
	CD

\end{enumerate}
}
%% memo:
\memoanswer{
（1）固定刻度读数为 0，可动刻度读数为 39.7，所测长度为 $0+39.7\times0.01=0.397\Umm$（0.395～0.399）。

（2）由记录数据根据欧姆定律可知金属丝的电阻 $R_{x}$ 约 $5\UO$。则有 $\frac{R_{x}}{R_{A}} = \frac{5}{0.1} = 50$，$\frac{R_{V}}{R_{x}} = \frac{3000}{5} = 600$，比较可知 $R_{x}$ 为小电阻，应该采用外接法测量误差小。由（3）知是用伏安特性曲线来测量电阻的，就要求电压电流从接近 0 开始调节，所以应该采用分压接法（甲）。

（3）注意连图时连线起点和终点在接线柱上并且不能交叉，结合（2）可知应该连接成外接分压接法（甲），那么在连线时断开开关且使 $R_{x}$ 两端的电压为 0。先连外接法（线 1）再连分压法（线 2 和 3），此时滑片 $P$ 必须置于变阻器的左端，如答图 3 所示。

（4）描绘出第 2、4、6 三个点后可见第 6 次测量数据的坐标点误差太大，舍去，然后作出 $U-I$ 图线，如答图 4 所示。其中第 4 次测量数据的坐标点在描绘出的 $U-I$ 图线上，有 $R_{x} = \frac{1.00}{0.220}\UO = 4.5\UO$（4.3～4.7）。

（5）根据电阻定律 $R_{x} = \rho \frac{l}{S}$，有 $\rho = \frac{R_{x} S}{l} = \frac{4.5 \times 3.14 \times 0.397^{2} \times 10^{-6}}{0.5}\UOm = 4.5 \times 10^{-6}\UOm$，估算出的金属丝电阻率是 C。

（6）用螺旋测微器测量金属丝直径时，由于读数引起的误差属于偶然误差；由于电流表和电压表内阻引起的误差属于系统误差。答案 CD。
}

\item
%% number: 22
%% paperName: 2012年北京理综第 22 题 16 分
%% typeId: 6
%% type: 计算题
%% chapter: 曲线运动
%% point: 平抛运动
%% method: 无
%% score: 16
%% degree: 650
%% duplicateId: 0
%% body:
如图所示，质量为 $m$ 的小物块在粗糙水平桌面上做直线运动，经距离 $l$ 后以速度 $v$ 飞离桌面，最终落在水平地面上。已知 $l=1.4\Um$，$v=3.0\Ums$，$m=0.10\Ukg$，物块与桌面间的动摩擦因数 $\mu=0.25$，桌面高 $h=0.45\Um$，不计空气阻力，重力加速度 $g$ 取 $10\Umsq$。求：
\begin{enumerate}
	\item
	小物块落地点距飞出点的水平距离 $s$；
	\item
	小物块落地时的动能 $E_{k}$；
	\item
	小物块的初速度大小 $v_{0}$。
\end{enumerate}
\onepicture[h=2.6cm, align=right]{figs/22.png}
%% answer: 
\jdanswer{
\begin{enumerate}
	\item
	$0.90\Um$

	\item
	$0.90\UJ$

	\item
	$4.0\Ums$

\end{enumerate}
}
%% memo:
\memoanswer{
（1）由平抛运动规律，有竖直方向 $h=\frac{1}{2}gt^{2}$，水平方向 $s = vt$，得水平距离 $s = \sqrt{\frac{2h}{g}} v = 0.90\Um$。

（2）由机械能守恒定律，动能 $E_{k}=\frac{1}{2}mv^{2}+mgh=0.90\UJ$。

（3）由动能定理，有 $-\mu mgl = \frac{1}{2}mv^{2} - \frac{1}{2}mv_{0}^{2}$，得初速度大小 $v_{0}=\sqrt{2\mu gl+v^{2}}=4.0\Ums$。
}

\item
%% number: 23
%% paperName: 2012年北京理综第 23 题 18 分
%% typeId: 6
%% type: 计算题
%% chapter: 功和能
%% point: 动能定理
%% method: 图象法
%% score: 18
%% degree: 450
%% duplicateId: 0
%% body:
摩天大楼中一部直通高层的客运电梯，行程超过百米．电梯的简化模型如图 1 所示。考虑安全、舒适、省时等因素，电梯的加速度 $a$ 是随时间 $t$ 变化的，已知电梯在 $t = 0$ 时由静止开始上升，$a-t$ 图像如图 2 所示。电梯总质量 $m = 2.0\times10^{3}\Ukg$。忽略一切阻力，重力加速度 $g$ 取 $10\Umsq$。
\begin{enumerate}
	\item
	求电梯在上升过程中受到的最大拉力 $F_{1}$ 和最小拉力 $F_{2}$；
	\item
	类比是一种常用的研究方法。对于直线运动，教科书中讲解了由 $v-t$ 图像求位移的方法。请你借鉴此方法，对比加速度和速度的定义，根据图 2 所示 $a-t$ 图像，求电梯在第 1 s 内的速度改变量 $\Delta v_{1}$ 和第 2 s 末的速率 $v_{2}$；
	\item
	求电梯以最大速率上升时，拉力做功的功率 $P$；再求在 0\textasciitilde{}11 s 时间内，拉力和重力对电梯所做的总功 $W$。
\end{enumerate}
\onepicture[h=3.4cm, align=right]{\includesvg{figs/23}}
%% answer: 
\jdanswer{
\begin{enumerate}
	\item
	$2.2\times10^{4}\UN$，$1.8\times10^{4}\UN$

	\item
	$0.50\Ums$，$1.5\Ums$

	\item
	$P=2.0\times10^{5}\UW$，$W=1.0\times10^{5}\UJ$

\end{enumerate}
}
%% memo:
\memoanswer{
（1）由牛顿第二定律，有 $F - mg = ma$。由 $a-t$ 图像可知，$F_{1}$ 和 $F_{2}$ 对应的加速度分别是 $a_{1}=1.0\Umsq$，$a_{2}=-1.0\Umsq$。

$F_{1}=m(g+a_{1})=2.0\times10^{3}\times(10+1.0)\UN=2.2\times10^{4}\UN$

$F_{2}=m(g+a_{2})=2.0\times10^{3}\times(10-1.0)\UN=1.8\times10^{4}\UN$

（2）类比可得，所求速度变化量等于第 1 s 内 $a-t$ 图线下的面积 $\Delta v_{1}=0.50\Ums$。同理可得 $\Delta v_{2}=v_{2}-v_{0}=1.5\Ums$，$v_{0}=0$，第 2 s 末的速率 $v_{2}=1.5\Ums$。

（3）由 $a-t$ 图像可知，11 s\textasciitilde{}30 s 内速率最大，其值等于 0\textasciitilde{}11 s 内 $a-t$ 图线下的面积，有 $v_{m}=10\Ums$。此时电梯做匀速运动，拉力 $F$ 等于重力 $mg$，所求功率 $P=Fv_{m}=mg\cdot v_{m}=2.0\times10^{3}\times10\times10\UW=2.0\times10^{5}\UW$。由动能定理，总功 $W = E_{k2} - E_{k1} = \frac{1}{2}mv_{m}^{2} = \frac{1}{2} \times 2.0 \times 10^{3} \times 10^{2}\UJ = 1.0 \times 10^{5}\UJ$。
}

\item
%% number: 24
%% paperName: 2012年北京理综第 24 题 20 分
%% typeId: 6
%% type: 计算题
%% chapter: 运动和力
%% point: 动量守恒定律
%% method: 无
%% score: 20
%% degree: 450
%% duplicateId: 0
%% body:
匀强电场的方向沿 $x$ 轴正方向，电场强度 $E$ 随 $x$ 的分布如图所示，图中 $E_{0}$ 和 $d$ 均为已知量。将带正电的质点 A 在 O 点由静止释放，A 离开电场足够远后，再将另一带正电的质点 B 放在 O 点也由静止释放。当 B 在电场中运动时，A、B 间的相互作用力及相互作用能均为零；B 离开电场后，A、B 间的相互作用视为静电作用。已知 A 的电荷量为 $Q$，A 和 B 的质量分别为 $m$ 和 $\frac{m}{4}$。不计重力。
\begin{enumerate}
	\item
	求 $A$ 在电场中的运动时间 $t$；
	\item
	若 $B$ 的电荷量为 $q=\frac{4}{9}Q$，求两质点相互作用能的最大值 $E_{pm}$；
	\item
	为使 $B$ 离开电场后不改变运动方向，求 $B$ 所带电荷量的最大值 $q_{m}$。
\end{enumerate}
\onepicture[h=2.6cm, align=right]{figs/24.png}
%% answer: 
\jdanswer{
\begin{enumerate}
	\item
	$\sqrt{\frac{2dm}{QE_{0}}}$

	\item
	$\frac{1}{45}QE_{0}d$

	\item
	$\frac{16}{9}Q$

\end{enumerate}
}
%% memo:
\memoanswer{
（1）由牛顿第二定律，A 在电场中运动的加速度 $a=\frac{F}{m}=\frac{QE_{0}}{m}$。A 在电场中做匀变速直线运动，$d=\frac{1}{2}at^{2}$，解得运动时间 $t=\sqrt{\frac{2d}{a}}=\sqrt{\frac{2dm}{QE_{0}}}$。

（2）设 A、B 离开电场时的速度分别为 $v_{A0}$、$v_{B0}$，由动能定理，有 $QE_{0}d=\frac{1}{2}mv_{A0}^{2}$，$qE_{0}d=\frac{1}{2}\frac{m}{4}v_{B0}^{2}$。A、B 相互作用过程中，动量和能量守恒。A、B 相互作用力为斥力，A 受的力与其运动方向相同，B 受的力与其运动方向相反，相互作用力对 A 做正功，对 B 做负功。A、B 靠近的过程中，B 的路程大于 A 的路程，由于相互作用力大小相等，相互作用力对 B 做功的绝对值大于对 A 做功的绝对值，因此相互作用力做功之和为负，相互作用能增加。所以，当 A、B 最接近时相互作用能最大，此时两者速度相同，设为 $v'$，有 $(m+\frac{m}{4})v'=mv_{A0}+\frac{m}{4}v_{B0}$，$E_{pm}=(\frac{1}{2}mv_{A0}^{2}+\frac{1}{2}\frac{m}{4}v_{B0}^{2})-\frac{1}{2}(m+\frac{m}{4})v'^{2}$。已知 $q=\frac{4}{9}Q$，解得相互作用能的最大值 $E_{pm}=\frac{1}{45}QE_{0}d$。

（3）考虑 A、B 在 $x > d$ 区间的运动，由动量守恒、能量守恒，且在初态和末态均无相互作用，有 $mv_{A} + \frac{m}{4}v_{B} = mv_{A0} + \frac{m}{4}v_{B0}$，$\frac{1}{2}m v_{A}^{2}+\frac{1}{2}\frac{m}{4}v_{B}^{2}=\frac{1}{2}m v_{A0}^{2}+\frac{1}{2}\frac{m}{4}v_{B0}^{2}$。解得 $v_{B}=-\frac{3}{5}v_{B0}+\frac{8}{5}v_{A0}$。因 B 不改变运动方向，故 $v_{B}=-\frac{3}{5}v_{B0}+\frac{8}{5}v_{A0}\geq0$，解得 $q\leq\frac{16}{9}Q$，即 B 所带电荷量的最大值 $q_{m} = \frac{16}{9}Q$。
}

\end{enumerate}

\end{document}
